\documentclass[]{article}

\usepackage{amsmath}
\usepackage{multirow}
\usepackage{natbib}
\usepackage{graphicx} 


\begin{document}

Standard models of cosmic inflation rely on a postulated inflaton scalar field and its potential to drive the early universe's expansion. We present an alternative framework where inflation is realized as a metastable graviton condensate sustained by a self-regulating feedback mechanism. In this model, the quasi-de Sitter geometry is maintained by a balance between quantum depletion and a backreaction pressure from an information "memory burden" stored in the condensate's Bogoliubov modes. Through a combination of linear stability analysis and numerical integration, we demonstrate that this feedback loop creates a robust dynamical attractor. We show that fluctuations of the condensate naturally source the primordial curvature perturbations, correctly predicting a nearly scale-invariant, Gaussian, and red-tilted spectrum consistent with cosmological observations. A key finding is an information-theoretic constraint, $N_e \cdot N_s \leq M_{Pl}^2/H^2$, that links the inflationary duration ($N_e$) to the number of particle species ($N_s$) and the Hubble scale ($H$). Our simulations map the viable "stability corridor" in the $(N_s, H)$ parameter space where sufficient inflation occurs before the condensate's information capacity is saturated, leading to a natural exit via quantum breaking. Furthermore, a sensitivity analysis reveals that the inflationary energy scale can be dynamically selected by the particle content; for a linear scaling of the memory burden, the observed scale $H/M_{Pl} \sim 10^{-5}$ is uniquely determined by a particle count of $N_s \sim 10^5$, a value motivated by Grand Unified Theories. This work establishes a self-consistent alternative to standard inflation, replacing the inflaton potential with the information dynamics of a graviton condensate.

\end{document}
                